\documentclass[uplatex,a4paper,12pt]{jsarticle}
\usepackage{amsmath,amssymb}
\usepackage{enumitem}

\begin{document}

\title{機械学習基礎 確認テスト}
\author{}
\date{}
\maketitle

\section*{注意}
\begin{itemize}
\item 試験時間は5分とする．
\item 全10問，各10点，合計100点とする．
\item 60点以上を合格とする．
\end{itemize}

\section*{問題}

\begin{enumerate}

\item 機械学習の説明として最も適切なものを1つ選べ．
\begin{enumerate}[label=(\alph*)]
\item 人間がすべての判断ルールを記述する手法
\item データ集合からパターンや判断ルールを自動で獲得する手法
\item 乱数だけを用いて答えを決める手法
\end{enumerate}

\item 教師あり学習で用いるデータとして最も適切なものを1つ選べ．
\begin{enumerate}[label=(\alph*)]
\item 入力データのみ
\item 入力と出力のペア
\item 報酬のみ
\end{enumerate}

\item 次のうち，回帰問題の例として最も適切なものを1つ選べ．
\begin{enumerate}[label=(\alph*)]
\item 画像が犬か猫かを判定する
\item 手書き数字を0から9に分類する
\item 株価を予測する
\end{enumerate}

\item 次のうち，分類問題の説明として正しいものを1つ選べ．
\begin{enumerate}[label=(\alph*)]
\item 入力から連続値を予測する問題
\item 入力から離散的なクラスを予測する問題
\item 入力データの次元を増やす問題
\end{enumerate}

\item 次の単回帰モデル
\[
y=\alpha+\beta x
\]
において，重みに相当するものを1つ選べ．
\begin{enumerate}[label=(\alph*)]
\item $\alpha$
\item $\beta$
\item $y$
\end{enumerate}

\item 次の説明に当てはまる語句を1つ選べ．

「学習に使用していない未知のデータに対する予測性能」
\begin{enumerate}[label=(\alph*)]
\item 汎化性能
\item 過学習
\item バッチサイズ
\end{enumerate}

\item 次のうち，過学習の説明として正しいものを1つ選べ．
\begin{enumerate}[label=(\alph*)]
\item 学習データに過剰に適合し，未知データで性能が低下する現象
\item 学習データがまったく使われない現象
\item モデルの予測が常に正解する現象
\end{enumerate}

\item 二値分類において，
\[
TP=30,\quad TN=50,\quad FP=10,\quad FN=10
\]
であるとき，正解率 Accuracy を求めよ．

\item ソフトマックス関数の出力として正しいものを1つ選べ．
\begin{enumerate}[label=(\alph*)]
\item 各クラスに属する確率
\item 入力画像の画素値そのもの
\item 学習率
\end{enumerate}

\item 勾配降下法における学習率 $\eta$ の説明として正しいものを1つ選べ．
\begin{enumerate}[label=(\alph*)]
\item パラメータの更新量を調整する値
\item データの正解ラベル
\item クラス数
\end{enumerate}

\end{enumerate}

\end{document}